\chapter{Dental Cleaning}

\section*{Overview}
Dental cleaning removes plaque, tartar (calculus), and stains from the teeth to prevent periodontal disease, gingivitis, and dental cavities.

\section*{Alternative Names}
Prophylaxis, Dental scaling

\section*{Principle}
Scaling instruments and ultrasonic devices remove plaque and calculus above and below the gumline, and polishing removes surface stains. Eliminating the bacterial biofilm and irritants prevents gingivitis, periodontitis, and caries.
\section*{History}
Dental hygiene emerged as a formal profession in 1913 when Alfred Fones opened the first school of dental hygiene in Connecticut.

\section*{Why It Is Done}
Ordered as a routine preventive procedure every 6 months to maintain oral health and screen for oral cancers.

\section*{Is This a Primary or Secondary Test or Procedure}
Primary preventive treatment for gingivitis and periodontal disease.

\section*{How It Is Performed}
A dental hygienist uses ultrasonic instruments and hand scalers to scrape tartar off teeth, polishes the teeth with a gritty paste, and performs flossing.

\section*{Preparation}
Inform the hygienist if you require antibiotic prophylaxis before dental work (e.g., for specific heart valve conditions).

\section*{Contraindications and Precautions}
None. Delay if the patient has severe acute oral infections or uncorrected coagulopathy.

\section*{Risks}
Minor gum bleeding, mild tooth sensitivity, or rare transient bacteremia.

\section*{Aftercare and Recovery}
Resume normal brushing and flossing. Avoid eating or drinking for 30 minutes if a fluoride treatment was applied.

\section*{Outcome and Follow-up}
A dentist examines your clean teeth and X-rays for cavities, bone loss, or oral lesions.

\section*{Expected Results}
Removal of all supragingival and subgingival plaque and calculus; pink, healthy gingiva; no active decay.

\section*{Follow-up}
Routine dental cleaning in 6 months; restorative work if cavities were identified.

\section*{When to Seek Medical Care}
Follow the procedure team's specific instructions. Seek urgent medical assessment for severe or worsening pain, uncontrolled bleeding, fever or signs of infection, trouble breathing, chest pain, fainting, new weakness or confusion, inability to urinate, or any rapidly worsening symptom. The appropriate warning signs depend on the procedure and the patient's medical condition.

\section*{References}
\begin{enumerate}
  \item MedlinePlus. Dental Cleaning. U.S. National Library of Medicine; 2025.
  \item Current Medical Diagnosis and Treatment. McGraw-Hill; 2025.
\end{enumerate}

\clearpage
